\scanpage{38}

\begin{center}-- 98 --\end{center}
\margin{n\textsuperscript{o} 316}

un automorphisme de $E$.

Si en outre $E$ est une algèbre (non nécessairement associative), et si $u$ est une dérivation (nilpotente) de $E$, alors $e^u$ est un automorphisme de \uline{l'algèbre} $E$. En effet, si $x,y\in E$, on a
\[
u^p(xy)=\sum_{r+s=p}\binom{p}{r}u^r(x)u^s(y)
\]
pour tout entier $p\geqslant0$ (formule de Leibniz). Il en résulte que :
\[
\begin{aligned}
e^u(xy)&=\sum_{p\geqslant0}\frac{1}{p!}u^p(xy)=\sum_{p\geqslant0}\sum_{r+s=p}\frac{u^r(x)}{r!}\frac{u^s(y)}{s!}\\
&=\sum_{r,s=0}^\infty\frac{u^r(x)}{r!}\frac{u^s(y)}{s!}=e^u(x)e^u(y)
\end{aligned}
\]
d'où notre assertion.

Soit maintenant $\mathfrak g$ une algèbre de Lie. Si $x$ appartient au radical nilpotent de $\mathfrak g$, la dérivation $\operatorname{ad}_{\mathfrak g}x$ de $\mathfrak g$ est nilpotente. On peut donc poser la définition suivante :

\uline{Définition} 5 -- \uline{On appelle automorphisme spécial de} $\mathfrak g$ \uline{un automorphisme de} $\mathfrak g$ \uline{de la forme} $e^{\operatorname{ad}x}$, \uline{où} $x$ e\uline{st dans le radical nilpotent de} $\mathfrak g$.

Il est clair qu'un automorphisme spécial laisse stable tout idéal de $\mathfrak g$.

\uline{Définition} 6 -- \uline{Soient} $\mathfrak g$ \uline{une algèbre de Lie}, $\unsure{\mathfrak r|r}$ \uline{son radical}. \uline{On appelle sous-algèbre de Levi de} $\mathfrak g$ \uline{toute sous-algèbre supplémentaire de} $\unsure{\mathfrak r|r}$ \uline{dans} $\mathfrak g$.

Une sous-algèbre de Levi est isomorphe à $\mathfrak g/\unsure{\mathfrak r|r}$, donc semi-simple. Comme une sous-algèbre semi-simple n'a que $0$ en commun avec $\unsure{\mathfrak r|r}$, toute sous-algèbre semi-simple $\unsure{\mathfrak h|h}$ telle que $\mathfrak g=\unsure{\mathfrak r|r}+\unsure{\mathfrak h|h}$ est une sous-algèbre de Levi ; par conséquent, l'image d'une sous-algèbre de Levi par un homomorphisme surjectif est une sous-algèbre de Levi.
\margin{37}
